\subsubsection{\texttt{fnval}}

\texttt{fnval}函数返回样条在一个点或一组点处的函数值。

\begin{itemize}
    \item \texttt{y = fnval(s,x)}返回样条函数 \texttt{s} 在 \texttt{x} 处的值，\texttt{x}可以为标量、一维向量，
            \texttt{x}不能在\texttt{s}的定义域之外。返回值\texttt{y}的值取决于\texttt{x}。
    \item \texttt{fnval(s,x)}同\texttt{fnval(x,s)}。
\end{itemize}


\subsubsection{\texttt{fn2fm}}

\begin{itemize}
    \item \texttt{ns = fn2fm(s,form)} 将样条函数 \texttt{s} 转换为 \texttt{form} 格式，返回转换后的样条函数 \texttt{ns} 。 
    \item \texttt{form} 由字符向量或字符串标量形式指定，形式的选择分别为\texttt{'B-'}、\texttt{'pp'}，分别为 \texttt{B} 格式、\texttt{pp} 格式。
    \item \texttt{B-form} 将函数描述为给定结序列的给定阶 \texttt{k} 的 \texttt{B} 样条的加权和，\texttt{pp} 格式用其局部多项式系数来描述一个函数。

\end{itemize}

如输入：
\begin{lstlisting}[language = matlab]
    knots = [0,0,0,0,1,2,2,2,2];
    x  = [0,1,1,1,2];
    y  = [2,0,1,2,-1];
    s  = spapi(knots, x, y);
    ns = fn2fm(s,'pp');
\end{lstlisting}

 输出为：
\begin{lstlisting}[language = matlab ]
      form: 'pp'
    breaks: [0 1 2]
     coefs: [2 * 4 double]
    pieces: 2
     order: 4
       dim: 1
\end{lstlisting}
